% Source: Wald GR, physical PDF p. 450 (printed p. 437).
% Coverage: Appendix C opening overview.
% Figures/tables/footnotes: none.
% Status: source-reviewed and compile-clean.
% Uncertainties: none.

This appendix deals with topics related to the maps induced on tensor fields by maps
between manifolds. As will be shown in Section~C.1, if we have a map,
\(\phi:M\to N\), between manifolds \(M\) and \(N\), we can use \(\phi\) to bring
upper-index tensor fields from \(M\) to \(N\) and lower-index tensor fields from
\(N\) to \(M\). If \(\phi\) is a diffeomorphism, all types of tensor fields can be
carried from \(M\) to \(N\) or from \(N\) to \(M\). An important special case of this
result occurs when \(\phi_t:M\to M\) is a one-parameter family of diffeomorphisms
generated by a vector field \(v^a\). We can compare a given tensor field with the new
tensor field that arises from the action of \(\phi_t\) for small \(t\). As will be shown
in Section~C.2, this gives rise to the notion of the Lie derivative with respect to
the vector field \(v^a\). Finally, a vector field which generates a one-parameter
group of isometries is called a Killing vector field. Using the general formulas for
Lie derivatives, an equation for Killing fields is easily obtained and some important
properties of them are derived in Section~C.3.
